\ifdefined\gkver\else\def\gkver{student}\fi
\documentclass[\gkver]{gaokaozhenti}
\begin{document}

\chapter{2018年全国理综Ⅲ}
%% paperType: 理综物理部分

\begin{enumerate}
\item
%% number: 14
%% paperName: 2018年全国理综Ⅲ第 14 题 6 分
%% typeId: 1
%% type: 单选题
%% chapter: 原子和原子核
%% point: 人工核反应
%% method: 无
%% score: 6
%% degree: 850
%% duplicateId: 0
%% body:
1934 年，约里奥-居里夫妇用 $\alpha$ 粒子轰击铝核 $\ce{^{27}_{13}Al}$，产生了第一个人工放射性核素 X：$\ce{^{4}_{2}He + ^{27}_{13}Al -> ^{1}_{0}n + X}$。X 的原子序数和质量数分别为\xzanswer{B}

\fourchoices[answer=B]
{15 和 28}
{15 和 30}
{16 和 30}
{17 和 31}

%% answer: B
%% memo:
\memoanswer{
核反应过程遵守质量数守恒和电荷数守恒，其中 $\alpha$ 粒子的质量数为 4、电荷数为 2，中子 n 的质量数为 1、电荷数为 0，故 X 的质量数为 $4+27-1=30$，原子序数即电荷数为 $2+13-0=15$，故 B 正确。
}

\item
%% number: 15
%% paperName: 2018年全国理综Ⅲ第 15 题 6 分
%% typeId: 1
%% type: 单选题
%% chapter: 曲线运动
%% point: 开普勒定律
%% method: 无
%% score: 6
%% degree: 750
%% duplicateId: 0
%% body:
为了探测引力波，“天琴计划”预计发射地球卫星 P，其轨道半径约为地球半径的 16 倍；另一地球卫星 Q 的轨道半径约为地球半径的 4 倍。P 与 Q 的周期之比约为\xzanswer{C}

\fourchoices[answer=C]
{$2:1$}
{$4:1$}
{$8:1$}
{$16:1$}

%% answer: C
%% memo:
\memoanswer{
地球对卫星的万有引力提供卫星做圆周运动的向心力，设地球和卫星的质量分别为 $M$、$m$，由牛顿第二定律有 $G\frac{Mm}{r^{2}}=m\left(\frac{2\pi}{T}\right)^{2}r$，可得卫星的周期为 $T=\sqrt{\frac{4\pi^{2}r^{3}}{GM}}$，即卫星 P 与 Q 的周期之比 $\frac{T_{1}}{T_{2}}=\sqrt{\frac{r_{1}^{3}}{r_{2}^{3}}}=8$，故 C 正确。
}

\item
%% number: 16
%% paperName: 2018年全国理综Ⅲ第 16 题 6 分
%% typeId: 1
%% type: 单选题
%% chapter: 交流电
%% point: 交流电
%% method: 无
%% score: 6
%% degree: 650
%% duplicateId: 0
%% body:
一电阻接到方波交流电源上，在一个周期内产生的热量为 $Q_{\text{方}}$；若该电阻接到正弦交变电源上，在一个周期内产生的热量为 $Q_{\text{正}}$。该电阻上电压的峰值为 $u_{0}$，周期为 $T$，如图所示。则 $Q_{\text{方}}:Q_{\text{正}}$ 等于\xzanswer{D}

\onepicture[w=8cm]{figs/16.png}

\fourchoices[answer=D]
{$1:\sqrt{2}$}
{$\sqrt{2}:1$}
{$1:2$}
{$2:1$}

%% answer: D
%% memo:
\memoanswer{
由题图可知，方波交流电电压的有效值为 $u_{0}$，正弦交流电电压的有效值为 $\frac{u_{0}}{\sqrt{2}}$。设该电阻的阻值为 $R$，则 $Q_{\text{方}}=\frac{u_{0}^{2}}{R}T$，$Q_{\text{正}}=\frac{\left(\frac{u_{0}}{\sqrt{2}}\right)^{2}}{R}T=\frac{u_{0}^{2}}{2R}T$，即 $Q_{\text{方}}:Q_{\text{正}}=2:1$，故 D 正确。
}

\item
%% number: 17
%% paperName: 2018年全国理综Ⅲ第 17 题 6 分
%% typeId: 1
%% type: 单选题
%% chapter: 曲线运动
%% point: 平抛运动
%% method: 无
%% score: 6
%% degree: 450
%% duplicateId: 0
%% body:
在一斜面顶端，将甲乙两个小球分别以 $v$ 和 $\frac{v}{2}$ 的速度沿同一方向水平抛出，两球都落在该斜面上。甲球落至斜面时的速率是乙球落至斜面时速率的\xzanswer{A}

\fourchoices[answer=A]
{2 倍}
{4 倍}
{6 倍}
{8 倍}

%% answer: A
%% memo:
\memoanswer{
设斜面的倾角为 $\theta$，甲球落在斜面上所用时间为 $t$，由平抛运动的规律有 $x=vt$、$y=\frac{1}{2}gt^{2}$，且 $\tan\theta=\frac{y}{x}$，联立以上各式可得甲球落在斜面上所用时间为 $t=\frac{2v\tan\theta}{g}$，竖直方向的分速度为 $v_{y}=gt=2v\tan\theta$，甲球落在斜面上时速率 $v_{1}=\sqrt{v^{2}+v_{y}^{2}}=v\sqrt{1+4\tan^{2}\theta}$，同理可得乙球落在斜面上时速率为 $v_{2}=\frac{v}{2}\sqrt{1+4\tan^{2}\theta}$，即甲球落至斜面时的速率是乙球落至斜面时速率的 2 倍，故 A 正确。
}

\item
%% number: 18
%% paperName: 2018年全国理综Ⅲ第 18 题 6 分
%% typeId: 2
%% type: 多选题
%% chapter: 直线运动
%% point: 运动图象
%% method: 无
%% score: 6
%% degree: 550
%% duplicateId: 0
%% body:
甲乙两车在同一平直公路上同向运动，甲做匀加速直线运动，乙做匀速直线运动。甲乙两车的位置 $x$ 随时间 $t$ 的变化如图所示。下列说法正确的是\xzanswer{CD}

\onepicture[w=4.5cm]{\includesvg{figs/18}}

\fourchoices[answer=CD]
{在 $t_{1}$ 时刻两车速度相等}
{从 0 到 $t_{1}$ 时间内，两车走过的路程相等}
{从 $t_{1}$ 到 $t_{2}$ 时间内，两车走过的路程相等}
{从 $t_{1}$ 到 $t_{2}$ 时间内的某时刻，两车速度相等}

%% answer: CD
%% memo:
\memoanswer{
位移-时间图像的斜率表示物体运动的速度，由图可知，在 $t_{1}$ 时刻两图像的斜率不相等，即 $t_{1}$ 时刻两车的速度不相等，故 A 错误；在 0 时刻，乙车的位置坐标为 0，甲车在乙车的前方，$t_{1}$ 时刻两车位置坐标相同，为 $x_{1}$，故从 0 到 $t_{1}$ 时间内，甲车走过的路程小于乙车走过的路程，故 B 错误；$t_{2}$ 时刻两车位置坐标相同，为 $x_{2}$，故从 $t_{1}$ 到 $t_{2}$ 时间内，两车走过的路程相等，故 C 正确；由图可知，在 $t_{1}$ 到 $t_{2}$ 时间内的某时刻，两图像的斜率相等，即两车的速度相等，故 D 正确。
}

\item
%% number: 19
%% paperName: 2018年全国理综Ⅲ第 19 题 6 分
%% typeId: 2
%% type: 多选题
%% chapter: 运动和力
%% point: 牛顿定律中的图象
%% method: 无
%% score: 6
%% degree: 450
%% duplicateId: 0
%% body:
地下矿井中的矿石装在矿车中，用电机通过竖井运送至地面。某竖井中矿车提升的速度大小 $v$ 随时间 $t$ 的变化关系如图所示，其中图线①②分别描述两次不同的提升过程，它们变速阶段加速度的大小都相同；两次提升的高度相同，提升的质量相等。不考虑摩擦阻力和空气阻力。对于第①次和第②次提升过程，\xzanswer{AC}

\onepicture[w=6.5cm]{figs/19.png}

\fourchoices[answer=AC]
{矿车上升所用的时间之比为 $4:5$}
{电机的最大牵引力之比为 $2:1$}
{电机输出的最大功率之比为 $2:1$}
{电机所做的功之比为 $4:5$}

%% answer: AC
%% memo:
\memoanswer{
速度-时间图像和时间轴围成的图形的面积表示位移，因两次提升的高度相同，则位移相等，有 $\frac{1}{2}v_{0}\cdot 2t_{0}=\frac{1}{2}\cdot\frac{1}{2}v_{0}\cdot(t_{2}-t_{0}+t_{2})$，可得图线②表示的矿车上升的时间为 $t_{2}=2.5t_{0}$，故两次提升所用时间的比值为 $2t_{0}:2.5t_{0}=4:5$，故 A 正确；矿车加速上升时电机的牵引力最大，设矿车上升的加速度为 $a$，两次提升的质量相等，由牛顿第二定律有 $F-mg=ma$，则电机的最大牵引力相等，故 B 错误；电机的输出功率 $P=Fv$，最大牵引力相等，最大速度之比为 $2:1$，所以电机输出的最大功率之比为 $2:1$，故 C 正确；矿车的初、末速度为零，由动能定理有 $W_{F}-mgh=0$，即电机做的功等于克服矿车及矿石的重力做功的大小，克服重力做功相等，即电机做的功相等，故 D 错误。
}

\item
%% number: 20
%% paperName: 2018年全国理综Ⅲ第 20 题 6 分
%% typeId: 2
%% type: 多选题
%% chapter: 电磁感应
%% point: 电磁感应中图象
%% method: 无
%% score: 6
%% degree: 450
%% duplicateId: 0
%% body:
如图 \ref{2018全国理综Ⅲ20a}，在同一平面内固定有一长直导线 PQ 和一导线框 R，R 在 PQ 的右侧。导线 PQ 中通有正弦交流电流 $i$，$i$ 的变化如图 \ref{2018全国理综Ⅲ20b} 所示，规定从 Q 到 P 为电流的正方向。导线框 R 中的感应电动势\xzanswer{AC}

\twopicture[numA=a, numB=b, labelA=2018全国理综Ⅲ20a, labelB=2018全国理综Ⅲ20b, hA=3.6cm, hB=3.6cm]
{figs/20a.png}
{figs/20b.png}

\fourchoices[answer=AC]
{在 $t=\frac{T}{4}$ 时为零}
{在 $t=\frac{T}{2}$ 时改变方向}
{在 $t=\frac{T}{2}$ 时最大，且沿顺时针方向}
{在 $t=T$ 时最大，且沿顺时针方向}

%% answer: AC
%% memo:
\memoanswer{
直导线中电流变化时，产生的磁场发生变化，则线框中产生感应电动势，由法拉第电磁感应定律，线框 R 中感应电动势的大小与直导线中电流变化率成正比，由图乙可知，在 $t=\frac{T}{4}$ 时刻，图线切线的斜率为零，即电流的变化率为零，则线框中的感应电动势为零，故 A 正确；由安培定则和楞次定律可以判断，在 $t=\frac{T}{4}$ 时刻，线框 R 中的感应电动势改变方向，$t=\frac{T}{2}$ 时刻方向不变且沿顺时针方向，且图像的切线斜率最大，即感应电动势最大，故 B 错误 C 正确；在 $t=T$ 时刻，感应电动势的方向沿逆时针方向，故 D 错误。
}

\item
%% number: 21
%% paperName: 2018年全国理综Ⅲ第 21 题 6 分
%% typeId: 2
%% type: 多选题
%% chapter: 电场
%% point: 带电粒子的直线运动
%% method: 无
%% score: 6
%% degree: 350
%% duplicateId: 0
%% body:
如图，一平行板电容器连接在直流电源上，电容器的极板水平，两微粒 a、b 所带电荷量大小相等、符号相反，使它们分别静止于电容器的上、下极板附近，与极板距离相等。现同时释放 a、b，它们由静止开始运动，在随后的某时刻 $t$，a、b 经过电容器两极板间下半区域的同一水平面，a、b 间的相互作用和重力可忽略。下列说法正确的是\xzanswer{BD}

\onepicture[w=5.0cm]{figs/21.png}

\fourchoices[answer=BD]
{a 的质量比 b 的大}
{在 $t$ 时刻，a 的动能比 b 的大}
{在 $t$ 时刻，a 和 b 的电势能相等}
{在 $t$ 时刻，a 和 b 的动量大小相等}

%% answer: BD
%% memo:
\memoanswer{
两微粒同时经过电容器两极板间下半区域的同一水平面，即微粒 a 的位移大于微粒 b 的位移，由 $x=\frac{1}{2}at^{2}$ 可知，微粒 a 的加速度较大，因 a、b 间的相互作用和重力可忽略，则 a、b 只受电场力的作用，由牛顿第二定律得 $qE=ma$，两微粒带电荷量大小相等，场强相同，则微粒 a 的质量较小，故 A 错误；

由动能定理知，微粒的动能等于电场力做的功，因微粒 a 的位移大，故电场力对微粒 a 做的功较多，即 a 的动能比 b 的大，故 B 正确；在 $t$ 时刻，a、b 在同一水平面上，即同一等势面上，电势相等，由电势能的表达式 $E_{p}=q\varphi$，两微粒的电性相反，故电势能大小不相等，故 C 错误；在 $0\sim t$ 时间内，两微粒运动时间相同、受到的电场力大小相等，由动量定理有 $Ft=p-0$，故 a 和 b 的动量大小相等，故 D 正确。
}

\item
%% number: 22
%% paperName: 2018年全国理综Ⅲ第 22 题 6 分
%% typeId: 4
%% type: 实验题
%% chapter: 直线运动
%% point: 自由落体运动
%% method: 无
%% score: 6
%% degree: 650
%% duplicateId: 0
%% body:
甲、乙两同学通过下面的实验测量人的反应时间。实验步骤如下：
\begin{steps}
	\item
	甲用两个手指轻轻捏住量程为 $L$ 的木尺上端，让木尺自然下垂。乙把手放在尺的下端（位置恰好处于 $L$ 刻度处，但未碰到尺），准备用手指夹住下落的尺。
	\item
	甲在不通知乙的情况下，突然松手，尺子下落；乙看到尺子下落后快速用手指夹住尺子。若夹住尺子的位置刻度为 $L_{1}$，重力加速度大小为 $g$，则乙的反应时间为\tkanswer[2.6cm]{$\sqrt{\frac{2(L-L_{1})}{g}}$}（用 $L$、$L_{1}$ 和 $g$ 表示）。
	\item
	已知当地的重力加速度大小为 $g=9.80\Umsq$，$L=30.0\Ucm$，$L_{1}=10.4\Ucm$，乙的反应时间为\tkanswer[1.4cm]{0.20}$\Us$。（结果保留 2 位有效数字）
	\item
	写出一条提高测量结果准确程度的建议：\tkanswer[4.5cm]{多次测量取平均值；初始时乙的手指尽可能接近尺子}。
\end{steps}

%% answer:
%% memo:
\memoanswer{
（2）木尺下落做自由落体运动，下落时间即为反应时间，由 $h=\frac{1}{2}gt^{2}$，木尺下落的高度 $h=L-L_{1}$，可得 $t=\sqrt{\frac{2h}{g}}=\sqrt{\frac{2(L-L_{1})}{g}}$；

（3）将数据代入上式可得乙的反应时间为 $0.20\Us$；

（4）为提高准确度，可以多次测量取平均值，或初始时乙的手指尽可能接近木尺等。
}

\item
%% number: 23
%% paperName: 2018年全国理综Ⅲ第 23 题 9 分
%% typeId: 4
%% type: 实验题
%% chapter: 电路
%% point: 电路实验
%% method: 无
%% score: 9
%% degree: 550
%% duplicateId: 0
%% body:
一课外实验小组用如图所示的电路测量某待测电阻 $R_{x}$ 的阻值，图中 $R_{0}$ 为标准定值电阻（$R_{0}=20.0\UO$）；V 可视为理想电压表。S$_{1}$ 为单刀开关，S$_{2}$ 为单刀双掷开关，$E$ 为电源，$R$ 为滑动变阻器。采用如下步骤完成实验：

\twopicture[numA=a, numB=b, labelA=2018全国理综Ⅲ23a, labelB=2018全国理综Ⅲ23b, hA=5.0cm, hB=5.0cm]
{figs/23a.png}
{figs/23b.png}

\begin{steps}
	\item
	按照实验原理线路图\ref{2018全国理综Ⅲ23a}，将图\ref{2018全国理综Ⅲ23b}中实物连线；
	\item
	将滑动变阻器滑动端置于适当位置，闭合 S$_{1}$；
	\item
	将开关 S$_{2}$ 掷于 1 端，改变滑动变阻器动端的位置，记下此时电压表 V 的示数 $U_{1}$；然后将 S$_{2}$ 掷于 2 端，记下此时电压表 V 的示数 $U_{2}$；
	\item
	待测电阻阻值的表达式 $R_{x}=$\tkanswer[3.2cm]{$\left(\frac{U_{2}}{U_{1}}-1\right)R_{0}$}（用 $R_{0}$、$U_{1}$、$U_{2}$ 表示）；
	\item
	重复上述步骤，得到如下数据：

	\begin{center}
	\begin{tblr}{|c|c|c|c|c|c|}
	\hline
	 & 1 & 2 & 3 & 4 & 5 \\
	\hline
	$U_{1}/\UV$ & 0.25 & 0.30 & 0.36 & 0.40 & 0.44 \\
	\hline
	$U_{2}/\UV$ & 0.86 & 1.03 & 1.22 & 1.36 & 1.49 \\
	\hline
	$\frac{U_{2}}{U_{1}}$ & 3.44 & 3.43 & 3.39 & 3.40 & 3.39 \\
	\hline
	\end{tblr}
	\end{center}
	\item
	利用上述 5 次测量所得 $\frac{U_{2}}{U_{1}}$ 的平均值，求得 $R_{x}=$\tkanswer[1.4cm]{48.2}$\UO$。（保留 1 位小数）
\end{steps}

%% answer:
%% memo:
\memoanswer{
（1）连接实物图时，导线不能交叉，要注意滑动变阻器的分压式接法，实物连线如图所示。

\onepicture[w=5.5cm]{figs/23_an.png}

（4）由题意可得，当开关 S$_{2}$ 掷于 1 端时，测得 $R_{0}$ 两端电压为 $U_{1}$，当开关 S$_{2}$ 掷于 2 端时，电路中 $R_{0}$、$R_{x}$ 串联，测得 $R_{0}$、$R_{x}$ 两端电压为 $U_{2}$，由欧姆定律得 $\frac{U_{1}}{R_{0}}=\frac{U_{2}-U_{1}}{R_{x}}$，解得 $R_{x}=\frac{(U_{2}-U_{1})R_{0}}{U_{1}}=\left(\frac{U_{2}}{U_{1}}-1\right)R_{0}$。

（6）将数据代入上式可得 $R_{x}=48.2\UO$。
}

\item
%% number: 24
%% paperName: 2018年全国理综Ⅲ第 24 题 12 分
%% typeId: 6
%% type: 计算题
%% chapter: 磁场
%% point: 带电粒子在磁场中的运动
%% method: 无
%% score: 12
%% degree: 550
%% duplicateId: 0
%% body:
如图，从离子源产生的甲、乙两种离子，由静止经加速电压 $U$ 加速后在纸面内水平向右运动，自 M 点垂直于磁场边界射入匀强磁场，磁场方向垂直于纸面向里，磁场左边界竖直。已知甲种离子射入磁场的速度大小为 $v_{1}$，并在磁场边界的 N 点射出；乙种离子在 MN 的中点射出；MN 长为 $l$。不计重力影响和离子间的相互作用。求：
\begin{enumerate}
	\item
	磁场的磁感应强度大小；
	\item
	甲、乙两种离子的比荷之比。
\end{enumerate}
\onepicture[w=7.5cm, align=right]{figs/24.png}

%% answer:
\jdanswer{
\begin{enumerate}
	\item
	$B=\frac{4U}{lv_{1}}$
	\item
	$1:4$
\end{enumerate}
}
%% memo:
\memoanswer{
（1）设甲种离子所带电荷量为 $q_{1}$、质量为 $m_{1}$，在磁场中做匀速圆周运动的半径为 $R_{1}$，磁场的磁感应强度大小为 $B$，由动能定理得 $q_{1}U=\frac{1}{2}m_{1}v_{1}^{2}$ ①，

由洛伦兹力公式和牛顿第二定律得 $q_{1}v_{1}B=m_{1}\frac{v_{1}^{2}}{R_{1}}$ ②，

由几何关系知 $2R_{1}=l$ ③，

由①②③式得 $B=\frac{4U}{lv_{1}}$ ④。

（2）设乙种离子所带电荷量为 $q_{2}$、质量为 $m_{2}$，射入磁场的速度为 $v_{2}$，在磁场中做匀速圆周运动的半径为 $R_{2}$，同理有 $q_{2}U=\frac{1}{2}m_{2}v_{2}^{2}$ ⑤，

$q_{2}v_{2}B=m_{2}\frac{v_{2}^{2}}{R_{2}}$ ⑥，

由题给条件 $2R_{2}=\frac{l}{2}$ ⑦，

由①②③⑤⑥⑦式得，甲、乙两种离子的比荷之比为 $\frac{q_{1}}{m_{1}}:\frac{q_{2}}{m_{2}}=1:4$ ⑧。
}

\item
%% number: 25
%% paperName: 2018年全国理综Ⅲ第 25 题 20 分
%% typeId: 6
%% type: 计算题
%% chapter: 曲线运动
%% point: 竖直平面圆周运动
%% method: 无
%% score: 20
%% degree: 250
%% duplicateId: 0
%% body:
如图，在竖直平面内，一半径为 $R$ 的光滑圆弧轨道 ABC 和水平轨道 PA 在 A 点相切。BC 为圆弧轨道的直径。O 为圆心，OA 和 OB 之间的夹角为 $\alpha$，$\sin\alpha=\frac{3}{5}$，一质量为 $m$ 的小球沿水平轨道向右运动，经 A 点沿圆弧轨道通过 C 点，落至水平轨道；在整个过程中，除受到重力及轨道作用力外，小球还一直受到一水平恒力的作用，已知小球在 C 点所受合力的方向指向圆心，且此时小球对轨道的压力恰好为零。重力加速度大小为 $g$。求：
\begin{enumerate}
	\item
	水平恒力的大小和小球到达 C 点时速度的大小；
	\item
	小球到达 A 点时动量的大小；
	\item
	小球从 C 点落至水平轨道所用的时间。
\end{enumerate}
\onepicture[w=7.5cm, align=right]{figs/25.png}

%% answer:
\jdanswer{
\begin{enumerate}
	\item
	$F_{0}=\frac{3}{4}mg$，$v=\frac{\sqrt{5gR}}{2}$
	\item
	$p=\frac{m\sqrt{23gR}}{2}$
	\item
	$t=\frac{3}{5}\sqrt{\frac{5R}{g}}$
\end{enumerate}
}
%% memo:
\memoanswer{
（1）设水平恒力的大小为 $F_{0}$，小球到达 C 点时所受合力的大小为 $F$。由力的合成法则得 $\frac{F_{0}}{mg}=\tan\alpha$ ①，$F^{2}=(mg)^{2}+F_{0}^{2}$ ②，

设小球到达 C 点时的速度大小为 $v$，由牛顿第二定律得 $F=m\frac{v^{2}}{R}$ ③，

由①②③式和题给数据得 $F_{0}=\frac{3}{4}mg$，$v=\frac{\sqrt{5gR}}{2}$。

（2）设小球到达 A 点的速度大小为 $v_{1}$，作 $CD\perp PA$，交 PA 于 D 点，由几何关系得 $DA=R\sin\alpha$ ④，$CD=R(1+\cos\alpha)$ ⑤，

由动能定理得 $-mg\cdot CD-F_{0}\cdot DA=\frac{1}{2}mv^{2}-\frac{1}{2}mv_{1}^{2}$ ⑥，

由④⑤⑥式和题给数据得，小球在 A 点的动量大小为 $p=mv_{1}=\frac{m\sqrt{23gR}}{2}$。

（3）小球离开 C 点后在竖直方向上作初速度不为零的匀加速运动，加速度大小为 $g$。设小球在竖直方向的初速度为 $v_{\perp}$，从 C 点落至水平轨道上所用时间为 $t$，由运动学公式得

\onepicture[w=6.5cm]{figs/25_an.png}

$v_{\perp}t+\frac{1}{2}gt^{2}=CD$ ⑦，$v_{\perp}=v\sin\alpha$ ⑧，

由⑤⑦⑧式和题给数据得 $t=\frac{3}{5}\sqrt{\frac{5R}{g}}$。
}

\item
%% number: 33
%% paperName: 2018年全国理综Ⅲ第 33 题 10 分
%% typeId: 6
%% type: 计算题
%% chapter: 气体性质
%% point: 气体连接体
%% method: 无
%% score: 10
%% degree: 550
%% duplicateId: 0
%% body:
（1）如图，一定量的理想气体从状态 a 变化到状态 b，其过程如 $p-V$ 图中从 a 到 b 的直线所示。在此过程中\xzanswer{BCD}

\onepicture[w=4.5cm]{figs/33a.png}

\fivechoices[answer=BCD]
{气体温度一直降低}
{气体内能一直增加}
{气体一直对外做功}
{气体一直从外界吸热}
{气体吸收的热量一直全部用于对外做功}

（2）在两端封闭、粗细均匀的 U 形细玻璃管内有一股水银柱，水银柱的两端各封闭有一段空气。当 U 形管两端竖直朝上时，左、右两边空气柱的长度分别为 $l_{1}=18.0\Ucm$ 和 $l_{2}=12.0\Ucm$，左边气体的压强为 $12.0\UcmHg$。现将 U 形管缓慢平放在水平桌面上，没有气体从管的一边通过水银逸入另一边。求 U 形管平放时两边空气柱的长度。在整个过程中，气体温度不变。

\onepicture[w=3.5cm, align=right]{figs/33b.png}

%% answer:
\jdanswer{$l_{1}'=22.5\Ucm$，$l_{2}'=7.5\Ucm$}
%% memo:
\memoanswer{
（1）由图可知，气体从状态 a 变化到状态 b 的过程中，气体压强一直增大、体积一直增大，由理想气体状态方程 $\frac{pV}{T}=C$，可知气体的温度一直升高，内能一直增加，故 A 错误 B 正确；气体体积一直增大，则气体一直对外做功，故 C 正确；由热力学第一定律 $\Delta U=Q+W$，可知气体一直从外界吸热，吸收的热量一部分用于增加气体内能，一部分用于对外做功，故 D 正确 E 错误。

（2）设 U 形管两端竖直朝上时，左、右两边气体的压强分别为 $p_{1}$ 和 $p_{2}$。U 形管水平放置时，两边气体的压强相等，设为 $p$，此时原左、右两边气柱的长度分别变为 $l_{1}'$ 和 $l_{2}'$。由力的平衡条件得 $p_{1}=p_{2}+\rho g(l_{1}-l_{2})$ ①，式中 $\rho$ 为水银的密度，$g$ 为重力加速度大小。

\onepicture[w=4cm]{figs/33b_an.png}

由玻意耳定律得 $p_{1}l_{1}=pl_{1}'$ ②，$p_{2}l_{2}=pl_{2}'$ ③，两边气柱长度的变化量大小相等，即 $l_{1}'-l_{1}=l_{2}-l_{2}'$ ④，由①②③④式和题给条件得 $l_{1}'=22.5\Ucm$，$l_{2}'=7.5\Ucm$。
}

\item
%% number: 34
%% paperName: 2018年全国理综Ⅲ第 34 题 10 分
%% typeId: 6
%% type: 计算题
%% chapter: 光学
%% point: 光的折射
%% method: 无
%% score: 10
%% degree: 550
%% duplicateId: 0
%% body:
（1）一列简谐横波沿 $x$ 轴正方向传播，在 $t=0$ 和 $t=0.20\Us$ 时的波形分别如图中实线和虚线所示。已知该波的周期 $T>0.20\Us$。下列说法正确的是\xzanswer{ACE}

\onepicture[w=8cm]{figs/34a.png}

\fivechoices[answer=ACE]
{波速为 $0.40\Ums$}
{波长为 $0.08\Um$}
{$x=0.08\Um$ 的质点在 $t=0.70\Us$ 时位于波谷}
{$x=0.08\Um$ 的质点在 $t=0.12\Us$ 时位于波谷}
{若此波传入另一介质中其波速变为 $0.80\Ums$，则它在该介质中的波长为 $0.32\Um$}

（2）如图，某同学在一张水平放置的白纸上画了一个小标记“$\cdot$”（图中 O 点），然后用横截面为等边三角形 ABC 的三棱镜压在这个标记上，小标记位于 AC 边上。D 位于 AB 边上，过 D 点做 AC 边的垂线交 AC 于 F。该同学在 D 点正上方向下顺着直线 DF 的方向观察。恰好可以看到小标记的像；过 O 点做 AB 边的垂线交直线 DF 于 E；$DE=2\Ucm$，$EF=1\Ucm$。求三棱镜的折射率。（不考虑光线在三棱镜中的反射）

\onepicture[w=5cm, align=right]{figs/34b.png}

%% answer:
\jdanswer{$n=\sqrt{3}$}
%% memo:
\memoanswer{
（1）由图可知，波长为 $\lambda=16\Ucm=0.16\Um$，由题意可知 $\left(n+\frac{1}{2}\right)T=0.20\Us$，$n=0,1,2,\cdots$，又因波的周期 $T>0.20\Us$，则波的周期为 $T=0.4\Us$，波速 $v=\frac{\lambda}{T}=0.40\Ums$，选项 A 正确，B 错误；由波的传播方向可以判断，$t=0$ 时，$x=0.08\Um$ 处的质点沿 $y$ 轴正方向运动，时间 $t=0.70\Us=T+\frac{3}{4}T$，故 $t=0.70\Us$ 时刻，该质点位于波谷，故 C 正确；$\frac{1}{4}T<0.12\Us<\frac{1}{2}T$，故 $t=0.12\Us$ 时，$x=0.08\Um$ 处的质点在平衡位置的上方，故 D 错误；波传入另一介质时周期不变，所以波长变为 $\lambda=vT=0.32\Um$，故 E 正确。

（2）过 D 点作 AB 边的法线 $NN'$，连接 OD，则 $\angle ODN=\alpha$ 为 O 点发出的光线在 D 点的入射角；设该光线在 D 点的折射角为 $\beta$，如图所示。

\onepicture[w=5.5cm]{figs/34b_an.png}

由折射定律得 $n\sin\alpha=\sin\beta$ ①，其中 $n$ 为三棱镜的折射率。由几何关系可知 $\beta=60^{\circ}$ ②，$\angle EOF=30^{\circ}$ ③，在 $\triangle OEF$ 中，有 $EF=OE\sin\angle EOF$ ④，由③④式和题给条件得 $OE=2\Ucm$ ⑤，由题给条件可知，$\triangle OED$ 为等腰三角形，有 $\alpha=30^{\circ}$ ⑥，由①②⑥式得 $n=\sqrt{3}$ ⑦。
}

\end{enumerate}

\end{document}
